\section{Related Work}
\label{sec:related}

From the perspective of where nonlinearity is introduced, existing nonlinear tensor decomposition methods can be roughly grouped into four categories.

The first category preserves the low-rank factor backbone while directly modifying the mapping from latent linear predictors to observed tensor entries, that is, it places nonlinearity at the entry-generation level. Early work in this direction extends classical decomposition models through generalized losses or generalized likelihoods. For example, generalized CP decomposition \cite{hong2020generalized} preserves low-rank parameterization while replacing squared error with a more general observation loss, allowing the model to adapt to binary data, count data, and robust modeling scenarios. Bayesian Poisson Tucker decomposition \cite{schein2016bayesian} introduces a Poisson observation model and a hierarchical Bayesian structure into the Tucker framework for modeling count-valued multiway event data. Going further, InfTucker \cite{xu2011infinite}, Distributed Flexible Nonlinear Tensor Factorization \cite{zhe2016distributed}, and Streaming Nonlinear Bayesian Tensor Decomposition \cite{pan2020streaming} introduce stronger nonparametric entry-wise nonlinear expressivity through kernel methods, Gaussian processes, or their sparse approximations. Overall, this line of work is the closest to ours, because it likewise seeks to strengthen the observation-generation mechanism without completely abandoning the low-rank tensor backbone. The cost, however, is usually heavier inference, more kernel or variational hyperparameters, and relatively weaker structural transparency.

The second category places nonlinearity in the generation, constraint, or organization of the factors themselves, with the core idea being to change how the factors are constructed. Such methods typically bend the traditional linear factor space into a constrained nonlinear representation space through deep priors, temporal networks, or geometric manifold assumptions. For example, DeepTensor \cite{saragadam2024deeptensor} generates low-rank factors by deep networks and uses the implicit regularization of deep priors to compensate for the limited expressivity of classical low-rank models on complex signal structures. Neural Tensor Factorization \cite{wu2019neural} and DeepCP \cite{liu2017deepcp} combine embeddings, temporal dependency, and neural generation mechanisms to handle temporal interactions and noisy observations. Neural Additive Tensor Decomposition \cite{ahn2024neural} seeks a compromise between expressivity and interpretability by applying neural-network transformations to latent components so as to strengthen sparse tensor modeling. Other work emphasizes geometric structure more strongly, such as tensor decomposition with nonlinear manifold modeling for 3D head-pose estimation \cite{derkach2019tensor}, where explicit manifold constraints on pose factors improve generalization and interpretability. Compared with entry-level nonlinearity, this category usually has greater modeling freedom, but it also depends more heavily on network architecture design or prior selection, and its theoretical analysis and optimization stability are more complicated.

The third category goes further by networkizing the tensor decomposition operator itself, that is, by directly embedding nonlinearity into mode products, core contractions, or hierarchical decomposition structures, thereby forming deep tensor models. The representative idea is not to attach a nonlinear component outside a traditional low-rank decomposition, but to reorganize Tucker, TT, and related structures into multi-layer, stackable nonlinear transformation units. For instance, Stacked Tucker Decomposition \cite{xu2024stacked} inserts nonlinear activations after mode products and constructs a deep Tucker structure through hierarchical stacking. NMTucker \cite{varolgunes2023nmtucker} explicitly introduces nonlinear components into a Matryoshka-style Tucker framework to improve expressivity in highly sparse completion tasks. Nonlinear Tensor Train \cite{wang2021nonlinear} inserts activations between TT-core contractions for deep network compression and representation learning. Compared with the previous two categories, the advantage of this line is that it can approach the expressive power of deep networks while preserving parameter sharing in tensor structures, but at the cost of weakening the directness of classical low-rank models in parameter meaning, structural interpretation, and theoretical characterization.

The fourth category does not necessarily modify tensor decomposition itself directly. Instead, it treats tensors as compact carriers for representing external nonlinear structures, and then uses decomposition for efficient modeling and computation. Such methods commonly arise in polynomial systems, Volterra systems, and nonlinear transform-domain representations, where the nonlinear structure is typically specified first by physical mechanisms, analytic expansions, or front-end transforms, and tensors are mainly responsible for compactly representing higher-order coefficients and accelerating numerical computation. Typical examples include TNMOR and STORM \cite{liu2015model,deng2016storm}, which first expand nonlinear circuits or dynamical systems via polynomial, Taylor, or Volterra series, then rearrange the coefficients of different orders into high-order tensors and apply CP or symmetric decomposition to construct reduced-order models that preserve higher-order nonlinearity. Another line of work, such as multidimensional nonlinear transform-domain tensor representation \cite{shi2025multidimensional}, builds data-adaptive nonlinear transform domains using convolutional networks and activation functions, and then imposes tensor low-rank constraints in that domain. Overall, these methods show that tensors can serve not only as nonlinear decomposers but also as compact representations of nonlinear structure. In most cases, however, the source of nonlinearity lies outside the tensor model rather than forming a controlled, explicit, and directly interpretable entry-level extension inside the classical Tucker backbone.

Compared with the above work, NTD-PL is closer in spirit to an extension of the entry-generation mechanism that preserves the Tucker backbone, but it still differs from approaches based on generalized likelihoods, kernel methods, Gaussian processes, and deep decoders. Our method does not rely on infinite-dimensional kernel mappings, heavy probabilistic inference, or deep generative networks. Instead, it introduces an explicit, finite-order, element-wise polynomial link on top of a shared Tucker latent tensor, so that nonlinear capability enters the model in a controlled way through only a few coefficients. This design allows NTD-PL to inherit the low-rank structure, parameter efficiency, and multiway interpretability of Tucker, while also describing entry-level residuals through an explicit degree expansion and yielding a clear nesting relation with Tucker together with the corresponding approximation analysis. Relative to the existing literature, the novelty of NTD-PL lies precisely here: it does not fully outsource nonlinearity to kernels, neural networks, or external polynomial systems, but instead constructs within the classical Tucker framework an explicit, analyzable, and degree-wise controllable nonlinear link mechanism, thereby striking a new balance among expressive power, structural transparency, and optimization practicality.
